% Appendix — SpeechPrint evaluation results
% Dependencies: booktabs, longtable (both in preamble). Nothing else.
% No tipa, no external macros, no fancy packages.
% IPA written only where the character is in standard T1/lmodern (a e i o u ae),
% otherwise plain orthographic or ASCII approximation.

\appendix

% ============================================================
\chapter{SpeechPrint Prosody Output: English Minimal-Pair Corpus}
\label{app:english}
% ============================================================

The English corpus (\texttt{audio\_2026-05-30\_19-01-35.wav}, 305.5\,s) contains
82 sentences across 21 prosodic categories recorded by the author.
No hand annotation is available for this recording.
The analysis was produced by the full SpeechPrint pipeline: WhisperX (large-v3,
English model) for transcription and forced alignment, and Librosa pYIN for F0
extraction. All five algorithmic optimisations are applied (MAD-based adaptive
threshold, utterance boundary reset, declination removal, nucleus edge trimming,
octave recovery). Prosody symbols are computed relative to within-utterance
neighbours with MAD-based adaptive thresholds (floor 0.5\,ST,
factor $0.35 \times \hat{\sigma}_\mathrm{MAD}$).

\medskip\noindent\textbf{Prosody symbol key:}

\begin{center}
\small
\begin{tabular}{ll}
\toprule
Symbol & Meaning \\
\midrule
\texttt{/}   & weakly rising pitch within the syllable \\
\texttt{//}  & strongly rising (${>}\ 2.5\times$ adaptive threshold) \\
\texttt{\textbackslash}  & weakly falling \\
\texttt{\textbackslash\textbackslash} & strongly falling \\
\texttt{{\={}}}  & high level (U+203E overline; above neighbouring syllables) \\
\texttt{\_}  & low level (below neighbouring syllables) \\
\texttt{*}   & most prominent syllable in the utterance \\
\bottomrule
\end{tabular}
\end{center}

\section*{Sections 1--4: Selected examples with F0 detail}

\begin{table}[h]
\centering
\caption{Selected minimal-pair sentences with per-syllable F0 (onset\,$|$\,mid\,$|$\,offset
in Hz) and SpeechPrint prosody symbol.
`Prom.\ syl' = syllable carrying \texttt{*}.}
\label{tab:english_selected}
\small
\begin{tabular}{llllll}
\toprule
ID & Category & Sentence & Prom.\ syl & F0 (Hz) & Symbol \\
\midrule
1a & contr.\ stress
   & \textit{Do you want tea or COFFEE?}
   & k-o (coffee)  & 203$|$188$|$184 & \texttt{*‾\textbackslash} \\
1b & contr.\ stress
   & \textit{Did you meet MARY or MARTHA?}
   & m-ei (Mary)   & 262$|$343$|$303 & \texttt{*‾//} \\
\midrule
2b & focus: subj.
   & \textit{MARY flew to Milan yesterday.}
   & m-e (Mary)    & 212$|$226$|$245 & \texttt{*‾//} \\
2d & focus: loc.
   & \textit{Mary flew to MILAN yesterday.}
   & l-an (Milan)  & 197$|$214$|$232 & \texttt{‾//} \\
2f & focus: time
   & \textit{Mary flew to Milan YESTERDAY.}
   & yes-          & 212$|$219$|$227 & \texttt{*‾//} \\
2g & focus: manner
   & \textit{How did Mary go to Milan?}
   & did           & 220$|$215$|$213 & \texttt{*‾\textbackslash} \\
\midrule
3a & stmt./q.
   & \textit{Mary has arrived already.} [stmt]
   & m-e (Mary)    & 202$|$203$|$213 & \texttt{*‾} \\
3b & stmt./q.
   & \textit{Mary has arrived already?} [q]
   & has           & 233$|$252$|$287 & \texttt{*‾//} \\
3c & stmt./q.
   & \textit{Mary has arrived ALREADY?} [incred.]
   & m-e (Mary)    & 203$|$205$|$216 & \texttt{*‾} \\
\midrule
4b & lex.\ stress
   & \textit{I need a new \textbf{PER}mit.} [noun]
   & per-          & 178$|$174$|$171 & \texttt{*‾} \\
4a & lex.\ stress
   & \textit{I need to per\textbf{MIT} this.} [verb]
   & -mit          & 185$|$190$|$198 & \texttt{*‾/} \\
4c & lex.\ stress
   & \textit{The \textbf{CON}tract was signed.} [noun]
   & con-          & 207$|$227$|$233 & \texttt{*‾//} \\
\bottomrule
\end{tabular}
\end{table}

\noindent
The statement/question pair (3a/3b) is a clear test case:
the declarative receives a level or slightly falling nucleus (\texttt{*‾}),
while the interrogative accent on \textit{has} shows a strongly rising contour
(\texttt{*‾//}, 233\,$\to$\,287\,Hz), consistent with a final rise in English
polar questions.
In the lexical-stress pair (4b/4a) the \texttt{*} marker shifts correctly
between the first and second syllable, confirming that forced alignment captures
the acoustic difference.

\section*{Sections 5--21: Summary table}

\begin{longtable}{lll}
\caption{Prosody symbols for all 21 corpus categories (one representative sentence
per category where multiple exist). Full per-syllable data available in
\texttt{out/appendix/prosody\_minimal\_pairs\_readable.txt}.} \\
\toprule
ID & Sentence (abbreviated) & Symbol \\
\midrule
\endfirsthead
\toprule
ID & Sentence & Symbol \\
\midrule
\endhead
\bottomrule
\endfoot
5a  & \textit{John bought a new CAR yesterday.} [broad focus]    & \texttt{*‾} \\
5f  & \textit{JOHN bought a new car yesterday.} [narrow focus]   & \texttt{*‾} \\
5i  & \textit{JOHN bought it.} [contrastive topic]               & \texttt{*‾} \\
6a  & \textit{ONLY JOHN called Maria.}                           & \texttt{*‾\textbackslash} \\
6b  & \textit{John only CALLED Maria.}                           & \texttt{*‾} \\
6c  & \textit{John called only MARIA.}                           & \texttt{*‾} \\
6d  & \textit{EVEN JOHN passed the test.}                        & \texttt{*‾\textbackslash} \\
7a  & \textit{You finished the report.} [statement]              & \texttt{*‾} \\
7b  & \textit{You finished the report?} [question]               & \texttt{*‾} \\
7c  & \textit{You finished the report, DIDN'T you.}              & \texttt{*‾\textbackslash} \\
7d  & \textit{You finished the report, didn't YOU?}              & \texttt{*‾} \\
8a  & \textit{Who are you meeting?} [wh-question]                & \texttt{*‾} \\
8b  & \textit{You're meeting WHO?} [echo question]               & \texttt{*‾} \\
9a  & \textit{Do you want tea or COFFEE?} [yes/no]               & \texttt{*‾\textbackslash} \\
9b  & \textit{Do you want TEA or COFFEE?} [alternative]          & \texttt{*‾//} \\
9c  & \textit{Do you want TEA, COFFEE, or\ldots?} [list]         & \texttt{*‾//} \\
10a & \textit{That was incredibly good.} [declarative]           & \texttt{*‾} \\
10b & \textit{That was incredibly GOOD!} [exclamative]           & \texttt{*‾} \\
11a & \textit{As for the weather, it's getting COLD.}            & \texttt{*‾\textbackslash} \\
11b & \textit{JOHN I like; MARY I don't.} [contrastive topic]    & \texttt{*‾} \\
12a & \textit{JOHN called Maria, too.}                           & \texttt{*‾} \\
12b & \textit{John called MARIA, too.}                           & \texttt{*‾} \\
12c & \textit{John CALLED Maria, too.}                           & \texttt{*‾} \\
13b & \textit{She submitted it YESTERDAY.} [new information]     & \texttt{*‾\textbackslash} \\
13e & \textit{Emma submitted the draft YESTERDAY.}               & \texttt{*‾//} \\
14a & \textit{We need APPLES, PEARS, and BANANAS.}               & \texttt{*‾} \\
14b & \textit{We need APPLES, PEARS, and\ldots} [incomplete]     & \texttt{*‾\textbackslash} \\
15a & \textit{OLD men and women.}                                & \texttt{*‾\textbackslash\textbackslash} \\
15b & \textit{old MEN and WOMEN.}                                & \texttt{*‾} \\
16a & \textit{Are you comING?} [rising intonation]               & \texttt{*‾} \\
17a & \textit{Can you pass the SALT?} [request]                  & \texttt{*‾} \\
17b & \textit{Can you pass the SALT.} [command]                  & \texttt{*‾\textbackslash} \\
18a & \textit{My brother who LIVES in Boston called.} [restr.]   & \texttt{*‾} \\
18b & \textit{My BROTHER, who lives in BOSTON, CALLED.} [non-restr.] & \texttt{*‾\textbackslash} \\
19a & \textit{You met with the DIRECTOR?} [confirmation]         & \texttt{*‾} \\
19b & \textit{I met with the PRODUCER, not the DIRECTOR.}        & \texttt{*‾} \\
20a & \textit{John did NOT call Maria.}                          & \texttt{*‾//} \\
20b & \textit{John did not call MARIA; he called SARA.}          & \texttt{*‾} \\
21a & \textit{He's thirty.} [statement]                          & \texttt{*‾} \\
21b & \textit{He's THIRTY?} [surprise]                           & \texttt{*‾} \\
\end{longtable}


% ============================================================
\chapter{German GToBI Training Sentences: Human Annotation vs.\ SpeechPrint}
\label{app:german}
% ============================================================

Five GToBI-annotated German training sentences were processed with SpeechPrint.
The original TextGrids contain two hand-annotated tiers: \textit{Wort} (word
boundaries) and \textit{Ton} (GToBI nuclear tone and boundary tone annotations).
The final pipeline uses the human \textit{Wort} boundaries as the sole source of
word-level timing, with linguistically verified German IPA phone sequences
hardcoded per word and distributed proportionally within those boundaries.

\medskip\noindent
\textbf{A note on MFA.}
Montreal Forced Aligner (MFA) was evaluated as a phone-alignment source but
abandoned after discovering that it produced wrong transcriptions for three of
the five sentences even with the large-v3 Whisper transcript as input:
\textit{gelbe} became ``gerbe'', \textit{einige Melonen} became ``einigen mal
lohnt'', and \textit{er sang die Lieder} became ``hesangli lieder''.
In addition, MFA word-boundary timing deviated from the human \textit{Wort} tier
by 66--430\,ms across all five sentences, making it incompatible with the
human-annotation tiers. The hardcoded IPA approach eliminates both sources of
error.

\section*{GToBI annotations and expected pitch direction}

\begin{table}[h]
\centering
\small
\begin{tabular}{lll}
\toprule
Sentence & GToBI annotation & Expected direction \\
\midrule
\textit{eine gelbe Banane}       & L+H* L-\%  & rising \\
\textit{einige Melonen}          & H+L* L-\%  & falling \\
\textit{er sang die Lieder}      & H+!H* L-\% & falling (downstepped) \\
\textit{er will die Rosen haben} & L*+H       & rising \\
\textit{ich wohne in Bern}       & L*+H L-\%  & rising \\
\bottomrule
\end{tabular}
\end{table}

\section*{SpeechPrint prosody output (pYIN + 5 optimisations, hardcoded IPA)}

All five sentences use linguistically correct German IPA.
Table~\ref{tab:german} shows the per-syllable F0 (onset\,$\to$\,offset at the
vowel nucleus) and the SpeechPrint prosody symbol.

\begin{table}[h]
\centering
\caption{Per-syllable F0 and prosody symbols for the five German GToBI sentences.
F0 onset\,$\to$\,offset in Hz. Syllable labels from linguistically verified
German IPA (hardcoded, not espeak-ng).
\texttt{*} = most prominent syllable; other symbols as in
Appendix~\ref{app:english}.}
\label{tab:german}
\small
\setlength{\tabcolsep}{5pt}
\begin{tabular}{p{3.6cm} l l p{6.4cm}}
\toprule
Sentence & GToBI & Expected & Syllable: F0 onset$\to$offset (Hz) \quad symbol \\
\midrule
\textit{eine gelbe Banane}
  & L+H* L-\%
  & rising
  & eine-1: 205$\to$199\;\texttt{{\={}}} \quad
    eine-2: 199$\to$195\;\texttt{{\={}}} \quad
    gel: 192$\to$188\;\texttt{{\={}}} \quad
    ba: 193$\to$204\;\texttt{{\={}}} \quad
    \textbf{na}: 204$\to$261\;\texttt{//} \quad
    *na:\ 261$\to$182\;\texttt{\textbackslash\textbackslash} \\
\addlinespace[4pt]
\textit{einige Melonen}
  & H+L* L-\%
  & falling
  & ei: 115$\to$118\;\texttt{{\={}}} \quad
    ni: 118$\to$140\;\texttt{//} \quad
    ge: 140$\to$163\;\texttt{//} \quad
    *me: 159$\to$123\;\texttt{\textbackslash\textbackslash} \quad
    lo: 123$\to$104\;\texttt{\textbackslash} \quad
    nen: 104$\to$99\;\texttt{{\={}}} \\
\addlinespace[4pt]
\textit{er sang die Lieder}
  & H+!H* L-\%
  & falling
  & er: 103$\to$108\;\texttt{{\={}}} \quad
    *sang: 104$\to$165\;\texttt{//} \quad
    die: 166$\to$131\;\texttt{\textbackslash} \quad
    lie: 124$\to$101\;\texttt{\textbackslash} \quad
    der: 101$\to$103\;\texttt{{\={}}} \\
\addlinespace[4pt]
\textit{er will die Rosen haben}
  & L*+H
  & rising
  & er: 110$\to$108\;\texttt{{\={}}} \quad
    will: 109$\to$103\;\texttt{{\={}}} \quad
    die: 99$\to$104\;\texttt{{\={}}} \quad
    ro: 115$\to$157\;\texttt{//} \quad
    *sen: 157$\to$103\;\texttt{\textbackslash\textbackslash} \quad
    ha: 95$\to$90\;\texttt{{\={}}} \quad
    ben: 118$\to$100\;\texttt{\textbackslash} \\
\addlinespace[4pt]
\textit{ich wohne in Bern}
  & L*+H L-\%
  & rising
  & ich: 115$\to$104\;\texttt{{\={}}} \quad
    *woh: 113$\to$167\;\texttt{//} \quad
    ne: 167$\to$107\;\texttt{\textbackslash\textbackslash} \quad
    in: 102$\to$92\;\texttt{\textbackslash} \quad
    bern: 88$\to$107\;\texttt{/} \\
\bottomrule
\end{tabular}
\end{table}

\section*{Discussion}

\textbf{Rising accents (L+H*, L*+H).}
All three rising-accent sentences show a hat-shaped movement:
a strongly rising syllable (\texttt{//}) followed by a prominent falling peak
(\texttt{*\textbackslash\textbackslash}).
For \textit{Banane} (L+H*) the rise on the pre-nuclear syllable \textit{na}
(204\,$\to$\,261\,Hz, $+$4.3\,ST) and the fall from the peak
(261\,$\to$\,182\,Hz) are both captured correctly.
The phrase-final syllable \textit{bern} in \textit{ich wohne in Bern}
receives a slight rise (\texttt{/}, 88\,$\to$\,107\,Hz), consistent with a
continuation tone before the boundary.

\textbf{Falling accents (H+L*, H+!H*).}
For \textit{Melonen} (H+L*) the approach syllables \textit{ni} and \textit{ge}
carry \texttt{//} labels, and the nuclear syllable \textit{me} is marked
prominent with a strong fall (\texttt{*\textbackslash\textbackslash},
159\,$\to$\,123\,Hz), corresponding to L* low alignment.
For \textit{er sang die Lieder} (H+!H*) acoustic prominence falls on
\textit{sang} rather than on \textit{lieder}, because the largest pitch
excursion in the recording is the jump on \textit{sang} (104\,$\to$\,165\,Hz).
The Lieder syllables descend correctly.

\textbf{Pipeline note.}
For German GToBI, MFA and WhisperX are not used in the final pipeline. Word
boundaries come from the human \textit{Wort} tier; phone sequences are
hardcoded correct German IPA. The prosody layer (pYIN F0, five optimisations)
is entirely automatic and is the only component evaluated against the
human \textit{Ton} annotation.


% ============================================================
\chapter{Daakie (DoReCo): Italian and Spanish Model Comparison}
\label{app:daakie}
% ============================================================

The DoReCo corpus recording
\texttt{doreco\_port1286\_2017\_06\_30\_Jaklin.wav} (162.7\,s, 44100\,Hz, mono)
documents a fieldwork session in Daakie, an Oceanic language spoken on Malekula
Island, Vanuatu (ISO 639-3: \texttt{mtp}).
No automatic speech recognition model or pronunciation dictionary exists for Daakie.
SpeechPrint was run twice on the same recording using the nearest available
language models: Italian (\texttt{--language it}) and Spanish
(\texttt{--language es}), both with Whisper large-v3 via WhisperX.
Adaptive prosody thresholds: floor 0.5\,ST,
factor $0.35 \times \hat{\sigma}_\mathrm{MAD}$ (normalised median absolute deviation).

\section*{Transcription quality}

Neither model produces a correct orthographic transcription of Daakie.
Both generate plausible-sounding words from their own phonologies:
Italian yields forms such as \textit{s'occhini}, \textit{napia},
\textit{lumbino}; Spanish yields \textit{socine}, \textit{kolomdu},
\textit{san~wawa}.
The phoneme tier therefore reflects Italian or Spanish phonological categories
applied to Daakie acoustics, and should be treated as an acoustic scaffold
requiring manual correction for phonological work.

A phoneme-tier comparison against the DoReCo hand annotation is given in the body
of the thesis. The main finding is that the Italian model introduces segments absent
from Daakie (palatal nasal, lateral approximant, affricates, geminates) while
missing Daakie-specific features (labialized consonants, bilabial fricative,
pharyngeal /h/). The Spanish model introduces a different but similarly
mismatched inventory.

\section*{Coverage and prosody statistics}

\begin{table}[h]
\centering
\caption{Prosody statistics for Italian and Spanish mismatched-language conditions
on the Daakie DoReCo recording (162.7\,s).
v2 = updated adaptive thresholds (floor 0.5\,ST).
Old Italian v1 values (floor 3.0\,ST) in parentheses where they changed.}
\label{tab:daakie}
\small
\begin{tabular}{lrr}
\toprule
Metric & Italian (v2) & Spanish (v2) \\
\midrule
Words detected                    & 332           & 776 \\
Syllables                         & 572           & 918 \\
No-F0 syllables                   & 24\,(4.2\%)   & 96\,(10.5\%) \\
Mean F0 (Hz)                      & 194.6         & 194.2 \\
Static / level (\%)               & 76.2\,(v1: 90.3) & 67.0 \\
Weakly rising \texttt{/} (\%)     & 5.6           & 6.6 \\
Strongly rising \texttt{//} (\%)  & 4.0           & 3.6 \\
Weakly falling \texttt{\textbackslash} (\%)  & 7.0  & 9.8 \\
Strongly falling \texttt{\textbackslash\textbackslash} (\%) & 3.0 & 2.5 \\
Total rising (\%)                 & 9.6\,(v1: 4.0) & 10.2 \\
Total falling (\%)                & 10.0\,(v1: 2.8) & 12.3 \\
\bottomrule
\end{tabular}
\end{table}

\section*{Observations}

\textbf{Coverage.}
Spanish detects roughly 60\,\% more speech (918 vs.\ 572 syllables).
Higher recall means more of the recording receives a prosodic annotation
but also increases false-positive word detections in quieter sections,
reflected in the higher no-F0 rate (10.5\,\% vs.\ 4.2\,\%).

\textbf{Prosodic movement.}
Both conditions produce similar total rising+falling proportions
(Italian 19.6\,\%, Spanish 22.5\,\%), substantially higher than with the
original 3.0\,ST threshold (Italian v1: 6.8\,\%).
The updated thresholds detect approximately three times more pitch movement.
Strong-movement categories (\texttt{//} and
\texttt{\textbackslash\textbackslash}) make up roughly 6--7\,\% of syllables
combined, confirming that the double-symbol category is reserved for the
most prominent pitch excursions only.

\textbf{Prosody validity.}
Mean F0 is nearly identical across both conditions (194.6 vs.\ 194.2\,Hz),
confirming that the prosody symbols reflect the acoustic signal and are not
an artefact of the ASR language choice.

\textbf{Recommendation.}
The Italian/Spanish scaffolding documented here represents an earlier evaluation
stage. The final SpeechPrint pipeline for Daakie does not use ASR at all:
word and phone timing are taken directly from the DoReCo human annotations,
and CREPE is used for F0 extraction. That approach is described in
Section~\ref{sec:doreco_pipeline} and produces more reliable phoneme-tier
alignment than either mismatched ASR model can provide.
This appendix is retained as documentation of the cross-lingual phonological
projection phenomenon discussed in Section~\ref{sec:bias_projection}.


% ============================================================
\chapter{SpeechPrint: Installation and CLI Reference}
\label{app:speechprint_cli}
% ============================================================

\begin{center}
\fbox{\parbox{0.88\linewidth}{\centering\vspace{2mm}
\textbf{Navigate to the SpeechPrint repository:}\\[2mm]
\texttt{https://github.com/speechprint/SpeechPrint}\\[2mm]
\textbf{For Linux machines only.} macOS and Windows ports are planned.
\vspace{2mm}}}
\end{center}

\medskip

SpeechPrint is a prosody annotation environment and field-research toolkit.
It takes a WAV file as input and produces a Praat TextGrid as output,
combining forced alignment, phoneme extraction, F0 measurement across five
pitch trackers, and a symbolic prosody layer. The output is ready for Praat
without additional software installation.

\section*{Requirements}

\begin{center}
\small
\begin{tabular}{ll}
\toprule
Package & Installation \\
\midrule
\texttt{uv} & \texttt{pip install uv} \\
\texttt{librosa} & installed automatically \\
\texttt{torchcrepe} & installed automatically \\
\texttt{pesto-pitch} & installed automatically \\
\texttt{parselmouth} & installed automatically \\
\texttt{WhisperX} & configured by installer \\
MFA & configured by installer \\
\bottomrule
\end{tabular}
\end{center}

\section*{Setup}

\begin{verbatim}
git clone https://github.com/speechprint/SpeechPrint
cd SpeechPrint/linux
export SPEECHPRINT_ROOT="$PWD"
uv run python -m lib.main        # opens the launcher GUI
\end{verbatim}

On first run, select \textbf{Install SpeechPrint} to configure all
dependencies. Then select \textbf{New Project} to create an annotation
workspace.

\section*{CLI reference}

\subsection*{Annotate a single recording}

\begin{verbatim}
# English, pYIN tracker
python -m speechprint run \
    --wav recording.wav \
    --language en \
    --tracker pyin

# Endangered language with human-annotated TextGrid (DoReCo)
python -m speechprint run \
    --wav doreco.wav \
    --textgrid doreco.TextGrid \
    --tier-suffix "@TA" \
    --tracker crepe

# All five trackers, comparison TextGrid
python -m speechprint run \
    --wav recording.wav \
    --language en \
    --tracker pyin crepe pesto praat \
    --comparison
\end{verbatim}

\subsection*{Find a phonologically similar language}

\begin{verbatim}
# Suggest the closest supported ASR language for an endangered language
python -m speechprint suggest-language --iso "mtp"   # Daakie
python -m speechprint suggest-language --iso "cjp"   # Cabécar
\end{verbatim}

\subsection*{Batch processing}

\begin{verbatim}
# Annotate all WAV files in a folder
python -m speechprint batch \
    --input-dir corpora/recordings/ \
    --language en \
    --tracker pyin crepe \
    --comparison

# Export prosody tier to CSV
python -m speechprint export \
    --textgrid out/recording.TextGrid \
    --tier prosody \
    --format csv
\end{verbatim}

\subsection*{Evaluate against GToBI benchmark}

\begin{verbatim}
python -m speechprint evaluate-gtobi \
    --sentences-dir " five GToBI annotated sentences/" \
    --tracker pyin    # or crepe
\end{verbatim}

\section*{Output}

\begin{verbatim}
recording_name/
├── recording.TextGrid          ← open in Praat
├── words.csv
├── syllables.csv
├── phonemes.csv
├── prosody.csv
└── comparison/
    ├── prosody_pyin.TextGrid
    ├── prosody_crepe.TextGrid
    ├── prosody_pesto.TextGrid
    ├── prosody_praat.TextGrid
    └── tracker_comparison.csv
\end{verbatim}

\section*{Prosody symbol reference}

\begin{center}
\small
\begin{tabular}{ll}
\toprule
Symbol & Meaning \\
\midrule
\texttt{/}  & Weakly rising pitch within syllable \\
\texttt{//} & Strongly rising \\
\texttt{\textbackslash}  & Weakly falling \\
\texttt{\textbackslash\textbackslash} & Strongly falling \\
\texttt{{\={}}} & High level (U+203E; above neighbouring syllables) \\
\texttt{\_} & Low level (below neighbouring syllables) \\
\texttt{*}  & Prominent accent (louder + higher or longer) \\
\texttt{?}  & Unvoiced / insufficient F0 data \\
\bottomrule
\end{tabular}
\end{center}


% ============================================================
\chapter{Prosody Synthesis Toolkit: Articulatory Extensions}
\label{app:articulatory}
% ============================================================

\begin{center}
\fbox{\parbox{0.88\linewidth}{\centering\vspace{2mm}
A collection of experimental tools for prosody-driven articulatory synthesis.\\[2mm]
These tools connect the SpeechPrint symbolic prosody layer to
VocalTractLab articulatory synthesis, enabling closed-loop prosody
authoring: annotate intonation, then hear it resynthesised.\\[2mm]
\textbf{Prototype status.} Windows required for VocalTractLab synthesis.
\vspace{2mm}}}
\end{center}

\medskip

\begin{center}
\small
\begin{tabular}{ll}
\toprule
Package & Installation \\
\midrule
\texttt{tensortract2} & \texttt{pip install tensortract2} \\
\texttt{vocaltractlab-cython} & \texttt{pip install vocaltractlab-cython} \\
\texttt{pesto-pitch} & \texttt{pip install pesto-pitch} \\
\texttt{python-flucoma} & \texttt{https://github.com/jamesb93/python-flucoma} \\
\texttt{soundfile} & \texttt{pip install soundfile} \\
\texttt{scipy} & \texttt{pip install scipy} \\
\bottomrule
\end{tabular}
\end{center}

\section*{Setup}

\begin{verbatim}
# Create output folder from script directory
mkdir -p audio/output/
\end{verbatim}

\section*{Patches}

\subsection*{\texttt{prosody2tract.py}}

Convert a SpeechPrint prosody TextGrid tier to a VocalTractLab
tract-sequence file. Each prosody symbol is mapped to an F0 target
and duration, generating a source--target pair for the articulatory
trajectory.

\begin{verbatim}
# Basic conversion: prosody tier → VTL tract sequence
python prosody2tract.py recording.TextGrid

# Set F0 ceiling and floor from speaker range
python prosody2tract.py recording.TextGrid \
    --f0-floor 85 --f0-ceiling 300

# Use a specific movement generator for each pair
python prosody2tract.py recording.TextGrid \
    --generator sigmoid          # smooth S-curve
    --generator lorenz           # chaotic modulation
    --generator gendy            # Xenakis-style stochastic

# List available parameters and operations
python prosody2tract.py --list-params
python prosody2tract.py --list-generators
\end{verbatim}

\subsection*{\texttt{audio2tract.py}}

Audio-to-articulatory inversion. Converts a WAV file to a VTL
tract-sequence via TensorTract, optionally applying articulatory
parameter manipulations.

\begin{verbatim}
# Basic inversion (no manipulation)
python audio2tract.py my_audio.wav

# Raise pitch floor by scaling f0
python audio2tract.py my_audio.wav -m f0 set 150

# Increase tongue body constriction
python audio2tract.py my_audio.wav -m TCX multiply 1.5

# Smooth the subglottal pressure contour
python audio2tract.py my_audio.wav -m pressure smooth 10

# Multiple manipulations at once
python audio2tract.py my_audio.wav \
    -m TCX multiply 1.5 \
    -m TTY add 0.5 \
    -m f0 smooth 10

# Invert and write tract-sequence only (no synthesis)
python audio2tract.py my_audio.wav --no-synthesize
\end{verbatim}

\subsection*{\texttt{prosody\_morph.py} (Vox Prima / Vox Secunda)}

Morphs the prosodic contour of one speaker continuously into another.
Both recordings are inverted to articulatory trajectories; the morph
blends the F0 and vocal tract parameters at each frame.

\begin{verbatim}
# Morph speaker A prosody into speaker B (50% blend)
python prosody_morph.py \
    --prima speaker_a.wav \
    --secunda speaker_b.wav \
    --blend 0.5

# Animate the blend over the full duration
python prosody_morph.py \
    --prima speaker_a.wav \
    --secunda speaker_b.wav \
    --blend-curve sigmoid      # or lorenz, gendy
    --num-variations 5
\end{verbatim}

The output is a VTL tract-sequence file compatible with VocalTractLab~2.3
(\texttt{vocaltractlab.de}, Windows).

\section*{Movement generators}

Each source--target pair in the articulatory trajectory is shaped by a
movement generator that defines how parameters travel from source to target.

\begin{center}
\small
\begin{tabular}{lll}
\toprule
Generator & Category & Description \\
\midrule
\texttt{linear}      & Smooth    & Straight-line interpolation \\
\texttt{sigmoid}     & Smooth    & S-curve; slow at extremes, fast in middle \\
\texttt{ease\_in\_out} & Smooth  & Smoothstep cubic \\
\texttt{elastic}     & Smooth    & Overshoot-and-settle spring motion \\
\texttt{lorenz}      & Chaotic   & Lorenz attractor projected to one axis \\
\texttt{henon}       & Chaotic   & H\'enon map overlay \\
\texttt{duffing}     & Chaotic   & Duffing forced oscillator \\
\texttt{gendy}       & Stochastic & GENDY breakpoint synthesis (Xenakis, 1992) \\
\texttt{random\_walk} & Stochastic & Biased random walk arriving at target \\
\bottomrule
\end{tabular}
\end{center}

The \texttt{gendy} generator (adapted from Iannis Xenakis's GENDY
algorithm) is particularly relevant for prosody: it produces trajectories
where amplitude and duration at each breakpoint are independently updated
by random walks drawn from a Cauchy distribution, producing occasional
large jumps separated by small steps --- the characteristic behaviour of
emphatic and emotionally expressive speech intonation.

\section*{Parameter reference}

Tract parameters (indices 0--18):
\texttt{HX HY JX JA LP LD VS VO TCX TCY TTX TTY TBX TBY TRX TRY TS1 TS2 TS3}

Glottis parameters (indices 19--29):
\texttt{f0 pressure x\_bottom x\_top chink\_area lag rel\_amp double\_pulsing pulse\_skewness flutter aspiration\_strength}

\section*{Web-based TTS DAW (works now)}

\begin{center}
\fbox{\parbox{0.88\linewidth}{\centering\vspace{2mm}
\textbf{Pink Trombone Timeline Editor}\\[2mm]
Draw phoneme sequences on a timeline. Control articulation in real time.
Export audio directly from the browser.\\[2mm]
\texttt{https://zakaton.github.io/pink-trombone-editor/}\\[2mm]
A practical prototype of prosody-as-vector-graphics: the voice as an
editable object rather than a fixed recording.
\vspace{2mm}}}
\end{center}
